\begingroup
\setlength{\parindent}{0pt}
\setlength{\parskip}{0.40\baselineskip}
\setlength{\abovedisplayskip}{6.5pt plus 1.5pt minus 1.5pt}
\setlength{\belowdisplayskip}{6.5pt plus 1.5pt minus 1.5pt}
\setlength{\jot}{5pt}
\clubpenalty=10000
\widowpenalty=10000

\section*{How the Proof Works}

We start with a smooth, incompressible flow of finite energy and keep its
viscosity fixed and positive. The problem is to prove that the flow remains
smooth for all time. Its velocity \(u\) and pressure \(p\) satisfy the
three-dimensional Navier--Stokes equation:
\[
 \underbrace{\partial_tu}_{\text{local acceleration}}
 +\underbrace{(u\!\cdot\nabla)u}_{\text{nonlinear transport}}
 =\underbrace{-\nabla p}_{\text{pressure}}
 +\underbrace{\nu\Delta u}_{\text{viscous diffusion}},
 \qquad \nabla\!\cdot u=0.
\]
The hope is simple: viscosity smooths the fluid faster than nonlinear
transport can sharpen it.

Consider the case of a vortex within the fluid. Its vorticity
\(\omega=\nabla\times u\) satisfies
\[
 \partial_t\omega+(u\cdot\nabla)\omega
 =(\omega\cdot\nabla)u+\nu\Delta\omega.
\]
Strain within the fluid's flow can naturally stretch a vortex.
Incompressibility forces a tube of fluid particles along its core to narrow
as it lengthens.

{\centering
\includegraphics[
 width=0.96\linewidth,
 alt={Vortex stretching: axial strain sigma lengthens a material tube of length ell and radius r. Inward particle motion narrows it; v labels circumferential swirl and omega points along the vortex axis. Viscous diffusion spreads rotation across the material boundary on both sides.}
]{vortex-stretching-notation.pdf}\par}

For a thin, locally circular tube with length \(\ell\),
radius \(r\), and positive axial strain \(\sigma\),
\[
 \dot\ell=\sigma\ell>0,
 \qquad r\propto\ell^{-1/2}.
\]
If circulation \(\Gamma\) is retained through that narrowing, the mean
swirl speed \(v\) increases:
\[
 v=\frac{|\Gamma|}{2\pi r}.
\]
Viscosity diffuses momentum between neighbouring regions of fluid, reducing
velocity differences and spreading rotation through the surrounding fluid.
Over an elapsed time \(\tau\), diffusion reaches a distance of order
\[
 r_{\mathrm{diff}}\sim\sqrt{\nu\tau}.
\]
The vortex core can spread beyond the narrowing material tube, weakening
the circulation enclosed by it. Faster swirl depends on that circulation
remaining strong enough as the radius shrinks.

\needspace{10\baselineskip}
A narrower core gives viscosity less distance to act across, but the
growing velocity can also sharpen the flow more quickly. To compare these
effects at progressively smaller scales, we rescale the flow by
\(\lambda>1\):
\[
 \begin{aligned}
 u_\lambda(x,t)&=\lambda u(\lambda x,\lambda^2t),\\
 p_\lambda(x,t)&=\lambda^2p(\lambda x,\lambda^2t).
 \end{aligned}
\]
Lengths shrink by \(\lambda\), velocities grow by \(\lambda\), and times
shrink by \(\lambda^2\). Every term in the momentum equation gains a
factor \(\lambda^3\), leaving the rescaled flow governed by the original
equation at viscosity \(\nu\). Smaller scales alone give viscosity no
advantage over transport.

To measure the concentration, we need to count spatial variation as well
as the size of the velocity. The Sobolev index \(s\) sets how strongly
we count that variation: integer values count spatial derivatives in
\(L^2\), while fractional values weight frequency \(\xi\) by
\(|\xi|^s\). Velocity contributes \(\lambda\), differentiation
\(\lambda^s\), and the square root of the volume \(\lambda^{-3/2}\), giving
\[
 \|u_\lambda(\cdot,t)\|_{\dot H^s}
 =\lambda^{s-1/2}\|u(\cdot,\lambda^2t)\|_{\dot H^s}.
\]
Below \(s=\tfrac12\), concentration makes this norm smaller; above
\(s=\tfrac12\), it makes it larger. At exactly \(s=\tfrac12\), the
increased variation offsets the shrinking volume. This critical height
lets us compare concentration without the measurement changing simply
because we changed scale.

\needspace{7\baselineskip}
Energy sits below that height. For \(E=\tfrac12\|u\|_2^2\),
\[
 E[u_\lambda](t)=\lambda^{-1}E[u](\lambda^2t).
\]
The velocity grows while total energy falls: the faster motion occupies
less volume. An energy bound alone cannot exclude this concentration.

Nor does concentration necessarily run out of time. At width \(r\) and
velocity scale \(U\), transport acts over a time of order \(r/U\), and
viscosity over \(r^2/\nu\). Both shorten by \(\lambda^2\) under the
rescaling. If a flow could repeat these narrowing stages, their durations
would fit into a finite total:
\[
 \tau_j=\lambda^{-2j}\tau_0,
 \qquad
 \sum_{j=0}^{\infty}\tau_j
 =\frac{\tau_0}{1-\lambda^{-2}}<\infty.
\]
In this hypothetical Zeno cascade, infinitely many stages accumulate at
a finite time \(T_*\). There is no smallest width or final stage of
acceleration. The scaling allows this picture; the regularity problem
is whether the fluid can actually carry it out.

For a finite-time singularity, velocity would have to blow up. That requires
acceleration to blow up, then jerk, and so on, and so on. These changes in
velocity are governed by the Navier--Stokes equation. Differentiating it
shows what each further temporal order involves:
\[
 \begin{aligned}
 \partial_tu={}&\nu\Delta u-(u\cdot\nabla)u-\nabla p,\\
 \partial_t^2u={}&\nu\Delta\partial_tu
 -(\partial_tu\cdot\nabla)u\\
 &-(u\cdot\nabla)\partial_tu-\nabla\partial_tp,\\
 \partial_t^3u={}&\nu\Delta\partial_t^2u
 -(\partial_t^2u\cdot\nabla)u\\
 &-2(\partial_tu\cdot\nabla)\partial_tu
 -(u\cdot\nabla)\partial_t^2u
 -\nabla\partial_t^2p.
 \end{aligned}
\]
\needspace{10\baselineskip}
Each time derivative splits the transport product and applies a Laplacian
to the preceding response. Substituting the first equation into the second
exposes \(\Delta^2u\); the next substitution exposes \(\Delta^3u\).
Higher temporal orders reach further into the spatial derivatives, with
the nonlinear terms carried along:
\[
 \begin{aligned}
 \partial_t^{n+1}u={}&\nu\Delta\partial_t^nu\\
 &-\sum_{a=0}^{n}\binom na
 (\partial_t^au\cdot\nabla)\partial_t^{n-a}u
 -\nabla\partial_t^np.
 \end{aligned}
\]

At the critical height, this relation exposes the spatial regularity
behind those temporal demands. Write
\(E_s=\tfrac12\|\Lambda^su\|_2^2\), where
\(\Lambda=(-\Delta)^{1/2}\), and let \(H=E_{1/2}\).
Differentiating \(E_s\) and integrating the viscous term by parts moves
one spatial derivative onto each factor:
\[
 E_s'=-2\nu E_{s+1}+\mathcal P_s,
\]
with the transport and pressure contribution
\[
 \mathcal P_s=-\operatorname{Re}\langle\Lambda^su,
 \Lambda^s((u\cdot\nabla)u+\nabla p)\rangle_{L^2}.
\]
Starting at \(s=\tfrac12\), each further time derivative exposes the next
spatial energy:
\[
 \begin{aligned}
 H'&=-2\nu E_{3/2}+\mathcal P_{1/2},\\
 H''&=4\nu^2E_{5/2}-2\nu\mathcal P_{3/2}
        +\partial_t\mathcal P_{1/2}.
 \end{aligned}
\]
Continuing through \(H''',H^{(4)},\ldots\) reaches
\(E_{7/2},E_{9/2},\ldots\), together with the differentiated nonlinear
contributions.

\needspace{7\baselineskip}
Those spatial orders are enough for smoothness. Any chosen number \(k\)
of continuous spatial derivatives needs only a finite Sobolev height:
\[
 m>k+\tfrac32
 \quad\Longrightarrow\quad H^m(\mathbb R^3)\hookrightarrow C^k(\mathbb R^3).
\]
Every choice of \(k\) is covered by a sufficiently high half-integer order.
With finite energy, their intersection is
\[
 L^2\cap\bigcap_{n\ge0}\dot H^{n+1/2}
 =\bigcap_{m\ge0}H^m
 =H^\infty\hookrightarrow C^\infty.
\]
The temporal requirements of a singularity have led us into the spatial
requirements of smoothness. If these orders remain controlled together
through \(T_*\), their limiting velocity is smooth enough to start the
flow again.

If a singularity were to form, we would expect the opposite: smoothness at
every earlier instant, but an unbounded fixed Sobolev norm as \(t\) approaches
\(T_*\). The temporal responses would become unbounded too, even while total
energy remained finite. Having all finite derivatives before \(T_*\) would
not give us the control needed to carry them through it.

That control must come from the initial velocity
\(u_0\in H^\infty_\sigma(\mathbb R^3)\) under the equation.
Incompressibility fixes pressure through
\(-\Delta p=\partial_i\partial_j(u_i u_j)\), with its decay normalization.
Writing \(V_n=\partial_t^nu\) and \(Q_n=\partial_t^np\), the velocity
and pressure responses obey
\[
 \begin{aligned}
 V_{n+1}&=\nu\Delta V_n
 -\sum_{a=0}^n\binom na(V_a\cdot\nabla)V_{n-a}-\nabla Q_n,\\
 Q_n&=R_iR_j\sum_{a=0}^n\binom na V_{a,i}V_{n-a,j}.
 \end{aligned}
\]
\needspace{5\baselineskip}
Here \(R_i\) are the Riesz transforms and repeated spatial indices are
summed. Starting with \(V_0=u_0\) at \(t=0\), these relations determine
each subsequent response. Spatial differentiation fills out each temporal
row. The resulting collection of velocity and pressure derivatives is the
mixed jet generated by the datum.

We can retain any finite number of temporal rows and spatial orders.
Choosing a temporal depth \(N\) and spatial height \(M\) gives a rectangle
\(\mathcal R_{N,M}[u_0;T_*]\). The datum-generated rectangle theorem,
Theorem~\ref{thm:datum-rectangles}, controls its adjacent temporal rows
throughout the approach to \(T_*\):
\[
 \sup_{0<T<T_*}\sum_{n=0}^{N}\left(
 \|V_n\|_{L^2(0,T;H^{M+1})}^2
 +\|V_{n+1}\|_{L^2(0,T;H^{M-1})}^2
 \right)<\infty.
\]
The bound can depend on the orders \(N,M\), but remains finite as
\(T\) approaches \(T_*\). Since \(V_{n+1}=\partial_tV_n\), the
adjacent rows control both a response and its change in time. The spaces
\(H^{M+1}\) and \(H^{M-1}\) pair across \(H^M\); their time-integrable
product gives a continuous endpoint trace by the Hilbert-triple theorem:
\[
 V_n\in C([0,T_*];H^M).
\]

\needspace{11\baselineskip}
A larger rectangle retains the derivatives already present. On each
overlap, both rectangles contain derivatives of the original velocity and
pressure, so their entries and endpoint traces agree. We can increase
temporal depth and spatial height independently without losing that
agreement. Their compatible intersection gives
\[
 \begin{aligned}
 \mathcal I[u_0]
 &:=\bigl(\mathcal R_{N,M}[u_0;T_*]\bigr)_{N,M\ge0},\\
 u&\in\bigcap_{r,m\ge0}H^r(0,T_*;H^m(\mathbb R^3)).
 \end{aligned}
\]
Every finite request for derivatives fits inside one of these rectangles.
In particular, the velocity reaches \(T_*\) in every spatial Sobolev
space, and its traces agree as one limiting velocity:
\[
 \begin{aligned}
 u_*&=\lim_{t\uparrow T_*}u(t)\quad\text{in every }H^m,\\
 u_*&\in H^\infty_\sigma\hookrightarrow C^\infty.
 \end{aligned}
\]
Smooth local existence starts the flow again from \(u_*\), and uniqueness
joins it to the original solution. It continues beyond its supposed
maximal time. A finite \(T_*\) is impossible, so \(T_*=\infty\).

The constructed intersection now supplies finite estimates at any chosen
orders. For any finite time \(T\), temporal order \(r\), spatial order
\(k\), and \(2\le p\le\infty\), choose
\(m>k+\tfrac32-\tfrac3p\). Sobolev embedding gives
\[
 \partial_t^ru\in C([0,T];H^m)
 \hookrightarrow L^\infty(0,T;W^{k,p}).
\]
In particular, taking \(r=0\), \(k=1\), \(p=\infty\), and \(m=3\)
controls the total strain over any finite time:
\[
 \int_0^T\|\nabla u(t)\|_\infty\,dt
 \le T\sup_{0\le t\le T}\|\nabla u(t)\|_\infty<\infty.
\]
For the locally circular material tube at the start,
\(\dot r/r=-\sigma/2\) and \(\sigma\le\|\nabla u\|_\infty\). Hence
\[
 r(t)\ge r(0)\exp\!\left(
 -\frac12\int_0^t\|\nabla u(s)\|_\infty\,ds
 \right)>0.
\]
Only a finite amount of stretching can accumulate in a finite time.
The tube cannot be squeezed to zero radius there.

\endgroup
